\newpage
\section{Acta Número 05}
\label{sec:acta05}

\noindent \textbf{Datos de la Sesión}
\begin{itemize}[noitemsep]
    \item \textbf{Modalidad:} Virtual - Google Meet.
    \item \textbf{Hora de Inicio:} 18:00
    \item \textbf{Hora de Fin:} 20:00
\end{itemize}

\noindent \textbf{Orden del Día}
\begin{itemize}[noitemsep]
    \item \textbf{Asistencia.}
    \item \textbf{Revisión de las Epics Generadas:} Evaluación de los requerimientos de alto nivel y módulos principales del sistema.
    \item \textbf{Revisión de las Historias de Usuario (HU):} Análisis del desglose de las Epics, validación de criterios de aceptación y estimaciones preliminares.
    \item \textbf{Planificación de Repositorios:} Definición de la arquitectura de repositorios para Backend, Frontend y Documentación.
    \item \textbf{Revisión para los Repositorios de GitHub:} Establecimiento de políticas de ramas (branching), control de versiones y asignación de permisos.
    \item \textbf{Organización de Carpetas Físicas y Virtuales:} Estandarización de la estructura en Google Drive y manejo del archivo físico del equipo.
    \item \textbf{Firma y aprobación de los presentes.}
\end{itemize}

\noindent \textbf{Desarrolladores (Socios Fundadores)}
\begin{enumerate}[noitemsep]
    \item Pardo Pozo Juan Diego (Jefe de Proyecto).
    \item Arnez Jaimes Mauricio.
    \item Quispe Flores Camila Belen.
    \item Ariñez Bautista Jim Gabriel.
    \item Medina Rodríguez Mateo Jhoustin.
    \item Molina Maldonado Tomas Manuel.
    \item Zeballos Crespo Jonathan.
    \item Gutierrez Guzmán Angheline.
\end{enumerate}

\noindent \textbf{Fecha de la sesión:}\par
Jueves 19 de marzo de 2026

\newpage
\subsection{Desarrollo del Acta}

\noindent \textbf{Asistencia}\par
Se procedió a la toma de asistencia a la hora acordada, corroborando la presencia en sesión de los 8 socios fundadores de la grupo-empresa Byte Busters S.R.L. Al contar con la totalidad de los miembros, se estableció el quórum necesario para iniciar la toma de decisiones definitivas.

\vspace*{0.5cm}
\noindent\textbf{Revisión de las Epics Generadas}\par
Se llevó a cabo una lectura y análisis de las Epics definidas para el proyecto. El equipo verificó que estas características de alto nivel cubran íntegramente el alcance del "Sistema Generador de Portafolios Digitales", asegurando que los módulos principales estén claramente delimitados y alineados con los objetivos del cliente.

\vspace*{0.5cm}
\noindent \textbf{Revisión de las Historias de Usuario}\par
A partir de las Epics aprobadas, se revisó el desglose en Historias de Usuario (HU). Se prestó especial atención a la redacción orientada al valor del usuario, confirmando que cada HU cuente con sus respectivos Criterios de Aceptación para la fase de testing. Además, se realizó una verificación cruzada para asegurar que el tamaño de las historias sea manejable dentro de la duración de un Sprint (2 semanas).

\vspace*{0.5cm}
\noindent \textbf{Planificación y Revisión de Repositorios en GitHub}\par
Para garantizar un entorno de desarrollo ordenado, se aprobó la creación de tres repositorios independientes en GitHub: uno exclusivo para el \textbf{Backend}, otro para el \textbf{Frontend} y un tercero destinado al control de versiones de la \textbf{Documentación} (incluyendo archivos LaTeX y recursos gráficos). Posteriormente, se definieron las políticas de control de versiones, acordando el uso de un flujo de trabajo estándar (como GitFlow o GitHub Flow), estableciendo ramas principales (main/develop) y reglas estrictas para los Pull Requests y Code Reviews antes de cada integración.

\vspace*{0.5cm}
\noindent \textbf{Organización de Carpetas (Físicas y Google Drive)}\par
Con el objetivo de centralizar la información no alojada en código, se reestructuró la unidad compartida de Google Drive de la grupo-empresa. Se crearon directorios estandarizados para Administración, Desarrollo, Diseño y QA, definiendo nomenclaturas claras para los archivos. Paralelamente, se designó un responsable y un protocolo para el manejo de las carpetas físicas, asegurando el correcto archivado de contratos, actas firmadas y demás documentos formales requeridos por la consultora TIS.

\vspace*{0.5cm}
\noindent \textbf{Aprobación Oficial}\par
Habiendo estructurado la base técnica y organizativa del ciclo de desarrollo, los 8 socios fundadores aprobaron por unanimidad las decisiones tomadas respecto al Product Backlog y la gestión de la configuración (repositorios y directorios).

\newpage
\noindent \textbf{Firmas y aprobación de los integrantes}
\begin{center}
    \noindent
    \setlength{\tabcolsep}{0pt}
    
    \begin{tabular}{ccc}
        \espacioFirma{assets/img/firmas/mauricio.jpg}{Mauricio Jaimes Arnez}{13751221} & 
        \espacioFirma{assets/img/firmas/camila.jpg}{Camila Belen Quispe Flores}{13097244} &
        \espacioFirma{assets/img/firmas/jim.jpg}{Jim Gabriel Ariñez Bautista}{9517642} \\[1.0cm]
        
        \espacioFirma{assets/img/firmas/mateo.jpg}{Mateo Medina Rodríguez}{9453809} &
        \espacioFirma{assets/img/firmas/tomas.jpg}{Tomas Molina Maldonado}{8051701} &
        \espacioFirma{assets/img/firmas/jonathan.jpg}{Jonathan Zeballos Crespo}{13067494} \\[1.0cm]
    \end{tabular}

    \vspace{0.2cm} 
    \begin{tabular}{cc}
        \espacioFirma{assets/img/firmas/diego.jpg}{Juan Diego Pardo Pozo}{12747783} & 
        \espacioFirma{assets/img/firmas/angheline.jpg}{Angheline Gutierrez Guzmán}{13751815}
    \end{tabular}
\end{center}

\vspace{0.5cm}
\noindent \textbf{Fecha de última revisión:} 23 de marzo de 2026